\section{Introduction}\label{sec:introduction}

Can material inducements---foreign aid, loans, investment---buy a country's loyalty as expressed in diplomatic rhetoric? Kuziemko and Werker (2006) provided the canonical empirical strategy for this question: UN Security Council non-permanent seats, allocated by quasi-random regional rotation, attract additional US aid, which in turn should shift the recipient's foreign policy behavior toward US preferences. The subsequent literature has documented UNSC effects on IMF programs \citep{dreher2009global}, World Bank projects \citep{dreher2009development}, and UN voting alignment \citep{dreher2008does}. But whether aid-induced shifts in material incentives translate into \emph{expressed allegiance}---changes in how nations talk about the United States in their most visible annual diplomatic forum---remains an open question, and one that speaks directly to the political economy of influence.

This paper addresses that question using the UN General Debate corpus (Baturo et al.\ 2017; Jankin et al.\ 2024), a panel of US aid flows (Greenbook), and the UNSC membership instrument from Dreher, Sturm, and Vreeland (2009). We construct a semantic stance measure using large language model scoring (DeepSeek v4 Pro, version 3.1), validated against independent AI coding (Cohen's $\kappa = 0.65$, quadratic-weighted). We then estimate the first-stage relationship between UNSC membership and US aid, the reduced-form effect of UNSC membership on speech stance, and---conditional on passing a pre-registered causal gate---a 2SLS specification.

Our central finding is a null result with important methodological implications. The first stage fails decisively: across six fixed-effect specifications ranging from no controls to country--year FE, the \Fstat{} on UNSC membership never exceeds 1.1, far below the $F>10$ threshold required for reliable IV inference \citep{staiger1997instrumental}. The reduced form is consistently negative---UNSC members, if anything, express \emph{less} pro-US sentiment---with a point estimate of $-0.10$ ($p=0.075$) in the preferred country--year FE specification. Including GDP per capita, population, and trade openness as controls attenuates the estimate further ($-0.078$, $p=0.200$). An event study around the two-year UNSC term reveals a marginally significant positive coefficient two years before membership ($+0.132$, $p=0.009$), potentially reflecting election-year campaigning, but no systematic pattern during or after the term.

These results offer two contributions. First, we demonstrate that the UNSC--aid relationship that forms the foundation of the Kuziemko--Werker IV chain is absent from our data \emph{even when using their own region--year FE specification}. This is not merely a story of ``stricter identification kills the instrument''---the instrument was never alive. Second, we identify fundamental challenges in measuring expressed diplomatic allegiance through text: our v3.1 stance measure, despite passing AI-assisted validation, correlates only weakly with observable behavioral indicators like UN voting agreement ($r=0.085$) and ideal point estimates ($r=0.130$). The gap between \emph{what nations say} in diplomatic speeches and \emph{what they do} in formal votes is substantial and deserves further attention.

We report these null findings honestly, document every analytical choice in a reproducibility index, and make all code, data, and machine-readable results publicly available. The paper is structured as follows. Section~\ref{sec:literature} reviews the relevant literature. Section~\ref{sec:data} describes data sources and speech restoration. Section~\ref{sec:measurement} details stance measurement and validation. Section~\ref{sec:identification} presents the empirical strategy. Section~\ref{sec:results} reports results, and Section~\ref{sec:robustness} provides robustness checks. Section~\ref{sec:discussion} discusses implications and limitations.
